\chapter{CSS moderno e adattivo}

Nei capitoli 3 e 4 abbiamo visto come i selettori stabiliscono a quali elementi
si applica uno stile, come la cascata risolve le dichiarazioni in conflitto, il
box model e il flusso normale, il posizionamento, Flexbox e Grid, le media query
e i layout responsive. In questo capitolo estendiamo quei concetti con le
funzionalità più recenti del CSS: selettori più concisi, variabili, valori
calcolati, nuove unità per il viewport dei dispositivi mobili, query sui
contenitori, immagini che si adattano al dispositivo, preferenze dell'utente e
regole annidate.

\section{Selettori moderni}

I selettori CSS sono molto espressivi: combinando selettori semplici si
descrivono relazioni anche complesse fra gli elementi. Le liste di selettori
complessi, però, tendono a diventare ripetitive. I frammenti ripetuti allungano
il foglio di stile, che diventa più pesante da scaricare, e rendono le modifiche
più soggette a errori.

Strumenti come \term{Sass}, un preprocessore CSS che vedremo nel capitolo 19,
sono diventati popolari proprio perché offrivano modi comodi per organizzare e
semplificare i fogli di stile. Il CSS moderno ha incorporato direttamente nel
linguaggio diverse funzionalità di questo tipo.

\begin{esempio}{Un CSS ripetitivo}
Vogliamo dare lo stesso stile a tutti gli \tg{em} contenuti in un \tg{a}, in tre
diverse regioni della pagina. La regola a sinistra funziona, ma è ripetitiva. Se
poi volessimo limitarla agli \tg{em} dei link con classe \code{foo}, dovremmo
sostituire \code{a} con \code{a.foo} in tre punti del selettore.

\begin{affianco}[0.5][c]
\begin{lstlisting}[style=css, name=link.css]
header a em,
main a em,
footer a em {
  color: purple;
}
\end{lstlisting}
\risultato
\begin{lstlisting}[style=css, name=link-foo.css]
header a.foo em,
main a.foo em,
footer a.foo em {
  color: purple;
}
\end{lstlisting}
\wtnota{Tre occorrenze da modificare per un solo cambiamento.}
\end{affianco}
\end{esempio}

Le \term{pseudo-classi funzionali} aiutano a scrivere selettori più concisi. Si
chiamano così perché, come le funzioni, ricevono degli argomenti fra parentesi:
un elenco di selettori.

\begin{center}
\renewcommand{\arraystretch}{1.4}
\begin{tabular}{@{}l@{\hspace{1.2cm}}l@{}}
  \textbf{Pseudo-classe} & \textbf{A che cosa serve} \\
  \code{:is()}    & raggruppa selettori alternativi \\
  \code{:where()} & crea stili predefiniti facili da sovrascrivere \\
  \code{:not()}   & esclude gli elementi che corrispondono agli argomenti \\
  \code{:has()}   & seleziona un elemento in base ad altri elementi collegati \\
\end{tabular}
\end{center}

\subsection{:is()}

La pseudo-classe \code{:is()} (\textit{matches-any}) seleziona un elemento se
\textbf{almeno uno} dei selettori passati come argomento lo seleziona. È utile
quando di un selettore complesso cambia solo una parte: la parte che varia
finisce dentro \code{:is()}, quella comune si scrive una volta sola.

\begin{affianco}[0.48][c]
\begin{lstlisting}[style=css, name=prima.css]
header a em,
main a em,
footer a em {
  color: purple;
}

button:hover,
button:focus,
button:active {
  background: navy;
}
\end{lstlisting}
\risultato
\begin{lstlisting}[style=css, name=dopo.css]
:is(header, main, footer) a em {
  color: purple;
}

button:is(:hover, :focus, :active) {
  background: navy;
}
\end{lstlisting}
\wtnota{Le due versioni selezionano gli stessi elementi.}
\end{affianco}

Per limitare ora la prima regola ai link con classe \code{foo} basta una sola
modifica: \code{:is(header, main, footer) a.foo em}.

\subsection{:where()}

\code{:where()} seleziona gli elementi esattamente come \code{:is()}: un
elemento viene selezionato se almeno uno degli argomenti lo seleziona. La
differenza fondamentale è che \code{:where()} e \textbf{tutti i suoi argomenti}
contribuiscono alla cascata con specificità \textbf{zero}. Gli stili che ne
derivano sono quindi, di proposito, facili da sovrascrivere: è lo strumento
adatto per scrivere stili predefiniti, per esempio in una libreria, che chi la
usa potrà ridefinire con un selettore qualunque.

\begin{affianco}[0.44][c]
\begin{lstlisting}[style=css, name=where.css]
:where(main a #em) {
  color: red;
}

em {
  color: purple;
}
\end{lstlisting}
\begin{lstlisting}[style=html, name=where.html]
<body>
  <main>
    <a href="#">
      <em id="em">Emphasis</em> link
    </a>
  </main>
</body>
\end{lstlisting}
\risultato
\begin{browser}[where.html]
  \wtlink{\mbox{\textcolor{css-purple}{\textit{Emphasis}}} link}
\end{browser}
\begin{devtools}{Analisi pagina}
  \dtgruppo{Regole dell'elemento \texttt{<em id="em">}}
  \dtregola{em}{where.css:5}
  \dtdich{color}{purple}
  \dtchiudi
  \dtregola{:where(main a \#em)}{specificità (0,0,0) \quad where.css:1}
  \dtbarrata{color}{red}
  \dtchiudi
\end{devtools}
\wtnota{Vince \code{em}, con specificità (0,0,1),\\ anche se l'altro selettore
        contiene un id.}
\end{affianco}

Il selettore \code{main a \#em} avrebbe specificità $(1, 0, 2)$ e vincerebbe
senza discussioni; racchiuso in \code{:where()} vale $(0, 0, 0)$, e basta il
semplice selettore di tipo \code{em} per sovrascriverlo.

\subsection{:not()}

\code{:not()} seleziona gli elementi che \textbf{non} corrispondono a
\textbf{nessuno} dei suoi argomenti.

\begin{affianco}[0.5][c]
\begin{lstlisting}[style=css, name=not.css]
li:not(.selected, .active) {
  opacity: 0.5;
}
\end{lstlisting}
\begin{lstlisting}[style=html, name=lista.html]
<ol>
  <li>A</li>
  <li class="selected">B</li>
  <li class="active">C</li>
</ol>
\end{lstlisting}
\risultato
\begin{browser}[lista.html]
  \hspace*{1.2em}{\color{black!45}1.\enspace A}\par
  \hspace*{1.2em}2.\enspace B\par
  \hspace*{1.2em}3.\enspace C
\end{browser}
\wtnota{Solo A non ha né \code{.selected} né \code{.active}:\\ è l'unico
        elemento semitrasparente.}
\end{affianco}

\subsection{:has()}

\code{:has()} seleziona un elemento in base ai suoi \textbf{discendenti} o ad
altri elementi \textbf{collegati}. I suoi argomenti sono \term{selettori
relativi}: vengono valutati a partire dall'elemento stesso, per cui
\code{ol:has(.active)} indica gli \tg{ol} che hanno un discendente
\code{.active}, e \code{li:has(+ .selected)} i \tg{li} il cui fratello
\textbf{immediatamente successivo} ha la classe \code{selected}.

\begin{affianco}[0.55][c]
\begin{lstlisting}[style=css, name=has.css]
/* gli <ol> che contengono un .active */
ol:has(.active) {
  background-color: lightgray;
}

/* gli <li> seguiti subito da un .selected */
li:has(+ .selected) {
  color: red;
}
\end{lstlisting}
\risultato
\begin{browser}[lista.html]
  \noindent\colorbox{css-lightgray}{\parbox{\dimexpr\linewidth-2\fboxsep}{%
    \hspace*{1.2em}\textcolor{css-red}{1.\enspace A}\par
    \hspace*{1.2em}2.\enspace B\par
    \hspace*{1.2em}3.\enspace C}}
\end{browser}
\wtnota{L'\tg{ol} contiene C, che è \code{.active}: diventa grigio.\\ A è
        seguito da B, che è \code{.selected}: diventa rosso.}
\end{affianco}

Per la prima volta il CSS può selezionare un elemento in base a ciò che
\textbf{contiene} o a ciò che lo \textbf{segue}: per questo \code{:has()} viene
spesso chiamato \guillemotleft selettore del genitore\guillemotright.

\warning{I commenti CSS si scrivono solo con /* */}{
  Nelle slide i commenti di questo esempio sono scritti con \code{//}, come in
  JavaScript. In CSS non è un commento: il browser legge \code{//} e il resto
  della riga come l'inizio di una dichiarazione non valida, che si estende fino
  al primo punto e virgola, e la scarta insieme alla dichiarazione successiva
  (qui proprio \code{background-color}). L'unica sintassi per i commenti è
  \code{/* ... */}.}

\subsection{La specificità delle pseudo-classi funzionali}

Ricordiamo dal capitolo 3 che la specificità è una terna $(A, B, C)$: numero di
selettori di id; di classe, attributo e pseudo-classe; di tipo e
pseudo-elemento. Le pseudo-classi funzionali seguono regole proprie.

\begin{itemize}
  \item \textbf{Non} contano come pseudo-classi: non aggiungono nulla a $B$, a
        differenza di \code{:focus} o \code{:hover}.
  \item \code{:where()}, compreso tutto il suo contenuto, contribuisce sempre
        con specificità \textbf{zero}.
  \item Per \code{:is()}, \code{:not()} e \code{:has()} la specificità è quella
        del loro argomento \textbf{più specifico}.
\end{itemize}

\begin{affianco}[0.5][c]
\begin{lstlisting}[style=css, name=specificita.css]
li:is(.selected, #foo) {
  color: red;
}
\end{lstlisting}
\begin{lstlisting}[style=html, name=lista.html]
<ol>
  <li>A</li>
  <li class="selected">B</li>
  <li class="active">C</li>
</ol>
\end{lstlisting}
\risultato
\begin{browser}[lista.html]
  \hspace*{1.2em}1.\enspace A\par
  \hspace*{1.2em}\textcolor{css-red}{2.\enspace B}\par
  \hspace*{1.2em}3.\enspace C
\end{browser}
\begin{devtools}{Analisi pagina}
  \dtregola{li:is(.selected, \#foo)}{specificità (1,0,1)}
  \dtdich{color}{red}
  \dtchiudi
\end{devtools}
\wtnota{B è selezionato grazie a \code{.selected}, ma la regola\\ ha la
        specificità di \code{\#foo}, il suo argomento più specifico.}
\end{affianco}

La specificità non dipende da quale argomento seleziona davvero l'elemento: la
regola vale $(1, 0, 1)$ anche se nessun elemento della pagina ha l'id
\code{foo}. Ecco i selettori di questa sezione.

\begin{center}
\renewcommand{\arraystretch}{1.35}
\begin{tabular}{@{}l@{\hspace{1.2cm}}c@{\hspace{1.2cm}}l@{}}
  \textbf{Selettore} & \textbf{Specificità} & \textbf{Perché} \\
  \code{header a em}                         & $(0, 0, 3)$ & tre selettori di tipo \\
  \code{:is(header, main, footer) a em}      & $(0, 0, 3)$ & l'argomento più specifico è un tipo \\
  \code{:where(header, main, footer) a em}   & $(0, 0, 2)$ & \code{:where()} vale zero \\
  \code{:where(main a \#em)}                  & $(0, 0, 0)$ & tutto dentro \code{:where()} \\
  \code{li:not(.selected, .active)}          & $(0, 1, 1)$ & una classe e un tipo \\
  \code{li:has(+ .selected)}                 & $(0, 1, 1)$ & una classe e un tipo \\
  \code{li:is(.selected, \#foo)}              & $(1, 0, 1)$ & l'argomento più specifico è un id \\
  \code{button:is(:hover, :focus, :active)}  & $(0, 1, 1)$ & una pseudo-classe e un tipo \\
\end{tabular}
\end{center}

\section{Custom properties}

I valori ripetuti sono frequenti nei fogli di stile, e sono difficili da
mantenere.

\begin{lstlisting}[style=css, name=ripetuto.css]
h1     { color: #663399; }
a      { color: #663399; }
button { background: #663399; }
\end{lstlisting}

Che cosa succede quando il colore principale del sito cambia? Bisogna trovare e
sostituire ogni occorrenza, senza dimenticarne nessuna e senza cambiare per
errore lo stesso valore usato altrove con un altro significato.

\dfn{Custom property}{
  Le \term{custom property}, o \term{variabili CSS}, sono entità definite da chi
  scrive il foglio di stile. Rappresentano valori specifici, che si possono
  riutilizzare. Si definiscono con la sintassi delle custom property, cioè un
  nome che comincia con due trattini (\code{-{}-primary-color: \#663399}),
  oppure con la regola \code{@property}. Una volta definite, il loro valore si
  legge con la funzione \code{var()}.}

\begin{affianco}[0.5][c]
\begin{lstlisting}[style=css, name=variabili.css]
:root {
  --primary-color: #663399;
}

h1 {
  color: var(--primary-color);
}

a {
  color: var(--primary-color);
}

button {
  background: var(--primary-color);
}
\end{lstlisting}
\risultato
\begin{browser}[variabili.html]
  {\Lora\bfseries\LARGE\textcolor{css-rebeccapurple}{Titolo}}\par\vspace{6pt}
  \textcolor{css-rebeccapurple}{\uline{Un link}}\par\vspace{6pt}
  \wtbottone[css-rebeccapurple][black][black!75]{Pulsante}
\end{browser}
\wtnota{Per cambiare il colore di tutto il sito\\ basta modificare una riga.}
\end{affianco}

I nomi delle custom property distinguono maiuscole e minuscole:
\code{-{}-primary-color} e \code{-{}-Primary-Color} sono due proprietà diverse.

\subsection{Cascata ed ereditarietà}

Le custom property si comportano come le altre proprietà:

\begin{itemize}
  \item partecipano alla \textbf{cascata};
  \item sono normalmente \textbf{ereditate} dai discendenti;
  \item possono essere \textbf{sovrascritte}.
\end{itemize}

La pseudo-classe \code{:root} seleziona la radice dell'albero del documento (in
una pagina HTML, l'elemento \tg{html}) e si usa per definire custom property
\guillemotleft globali\guillemotright, che ogni elemento della pagina eredita.

\begin{affianco}[0.5][c]
\begin{lstlisting}[style=css, name=accent.css]
:root {
  --accent: #cfa2fb;
}

button {
  background: var(--accent);
}

.warning {
  --accent: orangered;
}
\end{lstlisting}
\begin{lstlisting}[style=html, name=accent.html]
<button>Normal button</button>
<button class="warning">Warning</button>
<div class="warning">
  <button>Inside a warning div</button>
</div>
\end{lstlisting}
\risultato
\begin{browser}[accent.html]
  \wtbottone[css-lilla][black][black!75]{Normal button}\enspace
  \wtbottone[css-orangered][black][black!75]{Warning}\par\vspace{4pt}
  \wtbottone[css-orangered][black][black!75]{Inside a warning div}
\end{browser}
\begin{immagine}[Da dove prende il suo valore ogni \code{var(-{}-accent)}.]
\begin{tikzpicture}[
    el/.style={mdom, anchor=west},
    v/.style={mnota, anchor=west},
    r/.style={mramo}]
  \node[el] (html) at (0,0) {<html>};
  \node[v] at ([xshift=4pt]html.east) {\code{-{}-accent: \#cfa2fb}};
  \node[el] (b1) at (0.6,-0.7) {<button>};
  \node[v] at ([xshift=4pt]b1.east) {eredita \textcolor{css-lilla!70!black}{\#cfa2fb}};
  \node[el] (b2) at (0.6,-1.4) {<button class="warning">};
  \node[v] at ([xshift=4pt]b2.east) {\textcolor{css-orangered}{orangered}, dichiarato qui};
  \node[el] (d) at (0.6,-2.1) {<div class="warning">};
  \node[v] at ([xshift=4pt]d.east) {\textcolor{css-orangered}{orangered}, dichiarato qui};
  \node[el] (b3) at (1.2,-2.8) {<button>};
  \node[v] at ([xshift=4pt]b3.east) {eredita \textcolor{css-orangered}{orangered}};
  \draw[r] (0.3,-0.22) |- (b1.west);
  \draw[r] (0.3,-0.22) |- (b2.west);
  \draw[r] (0.3,-0.22) |- (d.west);
  \draw[r] (0.9,-2.32) |- (b3.west);
\end{tikzpicture}
\end{immagine}
\end{affianco}

Il terzo pulsante non ha classi, ma è dentro un \tg{div} con classe
\code{warning}: la custom property \code{-{}-accent} viene ereditata dal
\tg{div}, e il pulsante diventa arancione come il secondo. (Per semplicità lo
schema omette \tg{body}, attraverso cui il valore di \tg{html} arriva ai
pulsanti.)

\subsection{Valori di riserva}

\code{var()} accetta un secondo argomento, il \term{valore di riserva}
(\textit{fallback}): se la custom property non è definita, si usa quello.

\begin{affianco}[0.5][c]
\begin{lstlisting}[style=css, name=notice.css]
.notice {
  color: var(--notice-color, darkred);
}

.teal-theme {
  --notice-color: teal;
}
\end{lstlisting}
\begin{lstlisting}[style=html, name=notice.html]
<p class="notice">Avviso</p>
<div class="teal-theme">
  <p class="notice">Avviso</p>
</div>
\end{lstlisting}
\risultato
\begin{browser}[notice.html]
  \textcolor{css-darkred}{Avviso}\par\vspace{6pt}
  \textcolor{css-teal}{Avviso}
\end{browser}
\wtnota{Nel primo paragrafo \code{-{}-notice-color} non è definita:\\ si usa
        il valore di riserva \code{darkred}.}
\end{affianco}

I valori di riserva sono utili quando un componente permette una
personalizzazione \textbf{facoltativa}: chi lo usa può definire
\code{-{}-notice-color}, e se non lo fa il componente ha comunque un aspetto
sensato.

\subsection{La regola @property}

La regola \code{@property} registra una custom property in modo esplicito:
ne dichiara il tipo di valore ammesso (\code{syntax}), se viene ereditata
(\code{inherits}) e il valore iniziale (\code{initial-value}, obbligatorio
salvo che il tipo sia \code{"*"}, cioè qualunque). Permette così, fra l'altro,
di personalizzare il comportamento dell'ereditarietà.

\begin{affianco}[0.5][c]
\begin{lstlisting}[style=css, name=property.css]
@property --accent {
  syntax: "<color>";
  inherits: false;
  initial-value: teal;
}

button {
  background: var(--accent);
}

.warning {
  --accent: orangered;
}
\end{lstlisting}
\risultato
\begin{browser}[accent.html]
  \wtbottone[css-teal][black][black!75]{Normal button}\enspace
  \wtbottone[css-orangered][black][black!75]{Warning}\par\vspace{4pt}
  \wtbottone[css-teal][black][black!75]{Inside a warning div}
\end{browser}
\wtnota{Stesso HTML di prima: ora il terzo pulsante\\ è di nuovo color
        \code{teal}.}
\end{affianco}

Con \code{inherits: false} il pulsante dentro il \tg{div} non eredita più il
valore \code{orangered} dal genitore: non dichiarando \code{-{}-accent} per
conto suo, prende il valore iniziale \code{teal}. Il secondo pulsante, che ha
lui stesso la classe \code{warning}, resta invece arancione.

\section{Valori calcolati}

\subsection{calc()}

La funzione \code{calc()} combina valori e unità \textbf{compatibili}, per
esempio una percentuale e dei pixel, e supporta le quattro operazioni
\code{+}, \code{-}, \code{*} e \code{/}.

\begin{affianco}[0.5][c]
\begin{lstlisting}[style=css, name=calc.css]
.container {
  background-color: purple;
  color: white;
  width: 200px;
}

.item {
  background-color: orangered;
  width: calc(100% - 20px);
  margin: auto;
}
\end{lstlisting}
\begin{lstlisting}[style=html, name=calc.html]
<div class="container">
  Container
  <div class="item">Item</div>
</div>
\end{lstlisting}
\risultato
\begin{browser}[calc.html]
\begin{tikzpicture}[font=\Lora\small,
    misura/.style={{Stealth[length=4pt]}-{Stealth[length=4pt]}, line width=0.6pt}]
  \fill[css-purple] (0,0) rectangle (4,1.0);
  \node[anchor=north west, text=white, inner sep=3pt] at (0,1.0) {Container};
  \fill[css-orangered] (0.2,0.02) rectangle (3.8,0.5);
  \node[anchor=west, text=white, inner sep=3pt] at (0.2,0.26) {Item};
  \draw[misura, draw=css-purple] (0,-0.25) -- (4,-0.25);
  \node[font=\Lato\scriptsize\bfseries, text=css-purple, anchor=north]
    at (2,-0.27) {200px};
  \draw[misura, draw=css-orangered] (0.2,-0.85) -- (3.8,-0.85);
  \node[font=\Lato\scriptsize\bfseries, text=css-orangered, anchor=north]
    at (2,-0.87) {100\% - 20px = 200px - 20px = 180px};
\end{tikzpicture}
\end{browser}
\wtnota{\code{100\%} è la larghezza del contenitore;\\ \code{margin: auto}
        lascia 10px per lato.}
\end{affianco}

\warning{Gli spazi attorno a + e -}{
  Dentro \code{calc()} gli operatori \code{+} e \code{-} vanno circondati da
  spazi: \code{calc(100\% - 20px)} è valido, \code{calc(100\%-20px)} no. Senza
  spazi il browser legge \code{-20px} come un numero negativo, non trova
  l'operatore e scarta l'intera dichiarazione. Per \code{*} e \code{/} gli spazi
  sono facoltativi.}

\subsection{min() e max()}

\code{min()} e \code{max()} restituiscono rispettivamente il \textbf{più
piccolo} e il \textbf{più grande} dei loro argomenti. Gli argomenti, separati
da virgole, possono essere più di due e possono contenere espressioni come
quelle di \code{calc()}.

\begin{affianco}[0.5][t]
\begin{lstlisting}[style=css, name=min.css]
.item {
  background-color: orangered;
  width: min(100px, calc(100% - 2rem));
  margin: auto;
}
\end{lstlisting}
L'elemento:
\begin{itemize}
  \item non è mai più largo di 100px;
  \item si restringe quando lo spazio disponibile è minore;
  \item conserva almeno 1rem di spazio per lato.
\end{itemize}
\risultato
\begin{lstlisting}[style=css, name=max.css]
.item {
  background-color: orangered;
  font-size: max(1rem, 10vw);
  margin: auto;
}
\end{lstlisting}
La dimensione del testo:
\begin{itemize}
  \item non è mai minore di 1rem;
  \item cresce nei viewport più larghi.
\end{itemize}
\end{affianco}

Dentro \code{min()} la funzione \code{calc()} si può anche omettere:
\code{min(100px, 100\% - 2rem)} è equivalente. I grafici mostrano come variano
i due valori, supponendo \code{1rem} pari a 16 pixel (il valore predefinito dei
browser).

\begin{figure}[H]
  \centering
  \begin{tikzpicture}[font=\Lato\scriptsize,
      asse/.style={-{Stealth[length=5pt]}, line width=0.7pt, draw=neutral!45!black},
      curva/.style={line width=1.4pt, draw=css-orangered},
      guida/.style={dashed, line width=0.5pt, draw=neutral!55!white}]
    % min(100px, 100% - 2rem): x in px del contenitore, 1px = 0.025cm
    \begin{scope}[x=0.025cm, y=0.025cm]
      \draw[asse] (0,0) -- (260,0);
      \node[anchor=north] at (130,-15) {larghezza del contenitore};
      \draw[asse] (0,0) -- (0,130) node[anchor=south, align=center]
        {larghezza dell'elemento};
      \draw[guida] (132,0) -- (132,100);
      \draw[guida] (0,100) -- (132,100);
      \node[anchor=east] at (0,100) {100px};
      \node[anchor=north] at (132,0) {132px};
      \node[anchor=north] at (32,0) {32px};
      \draw[curva] (32,0) -- (132,100) -- (250,100);
      \node[anchor=south east, text=css-orangered!80!black, rotate=45]
        at (88,58) {100\% - 2rem};
      \node[anchor=south, text=css-orangered!80!black] at (195,100) {100px};
    \end{scope}
    % max(1rem, 10vw): x in px del viewport, 1px = 0.0125cm; y: 1px = 0.075cm
    \begin{scope}[xshift=8.6cm, x=0.0125cm, y=0.075cm]
      \draw[asse] (0,0) -- (520,0);
      \node[anchor=north] at (260,-5) {larghezza del viewport};
      \draw[asse] (0,0) -- (0,46) node[anchor=south] {dimensione del testo};
      \draw[guida] (160,0) -- (160,16);
      \node[anchor=east] at (0,16) {16px};
      \node[anchor=north] at (160,0) {160px};
      \draw[curva] (0,16) -- (160,16) -- (420,42);
      \node[anchor=south, text=css-orangered!80!black] at (80,16) {1rem};
      \node[anchor=south east, text=css-orangered!80!black, rotate=21]
        at (330,33) {10vw};
    \end{scope}
  \end{tikzpicture}
  \caption{A sinistra \code{min(100px, 100\% - 2rem)}: segue la larghezza del
           contenitore fino a 100 pixel, poi si ferma. A destra
           \code{max(1rem, 10vw)}: resta a 16 pixel finché \code{10vw} non lo
           supera, cioè oltre un viewport di 160 pixel.}
\end{figure}

\subsection{clamp()}

\code{clamp(minimo, preferito, massimo)} vincola un valore fra un minimo e un
massimo: il valore cresce in modo fluido, ma non diventa mai più piccolo del
minimo né più grande del massimo. Equivale a
\codel{max(minimo, min(preferito, massimo))}.

\begin{affianco}[0.42][c]
\begin{lstlisting}[style=css, name=clamp.css]
h1 {
  font-size: clamp(
    2rem,   /* minimo */
    7vw,    /* preferito */
    3em     /* massimo */
  );
}
\end{lstlisting}
\risultato
\begin{immagine}[Con \code{1rem} = 16px e il testo del genitore a 16px.]
\begin{tikzpicture}[font=\Lato\scriptsize, x=0.0085cm, y=0.06cm,
    asse/.style={-{Stealth[length=5pt]}, line width=0.7pt, draw=neutral!45!black},
    curva/.style={line width=1.4pt, draw=css-orangered},
    guida/.style={dashed, line width=0.5pt, draw=neutral!55!white}]
  \fill[primary!6!white] (0,0) rectangle (457,56);
  \fill[green-gradient-1!40!white] (457,0) rectangle (686,56);
  \fill[primary!6!white] (686,0) rectangle (900,56);
  \draw[asse] (0,0) -- (920,0);
  \node[anchor=north] at (460,-6.5) {larghezza del viewport};
  \draw[asse] (0,0) -- (0,60) node[anchor=south] {dimensione di \code{h1}};
  \draw[guida] (457,0) -- (457,32) (686,0) -- (686,48);
  \node[anchor=east] at (0,32) {32px};
  \node[anchor=east] at (0,48) {48px};
  \node[anchor=north] at (457,0) {457px};
  \node[anchor=north] at (686,0) {686px};
  \draw[curva] (0,32) -- (457,32) -- (686,48) -- (900,48);
  \node[anchor=south, text=neutral!35!black] at (228,34) {minimo: 2rem};
  \node[anchor=south, text=neutral!35!black, align=center] at (571,50)
    {preferito:\\ 7vw};
  \node[anchor=south, text=neutral!35!black] at (793,50) {massimo: 3em};
\end{tikzpicture}
\end{immagine}
\end{affianco}

Il valore preferito \code{7vw} vale 32 pixel con un viewport largo 457 pixel e
48 pixel con uno largo 686: fra queste due larghezze il titolo cresce con la
finestra, prima e dopo resta fermo. Nella proprietà \code{font-size}, \code{em}
si riferisce alla dimensione del testo del \textbf{genitore}: con il valore
predefinito di 16 pixel, \code{3em} vale 48 pixel.

\section{L'altezza del viewport sui dispositivi mobili}

\code{100vh} rappresenta l'altezza del viewport. Nei browser desktop si basa
sulle dimensioni correnti della finestra; nei browser mobili la questione è più
delicata, perché i controlli del browser, come la barra degli indirizzi, possono
comparire e scomparire: per esempio, scorrendo la pagina verso il basso la barra
si nasconde e lo spazio per la pagina aumenta.

Il CSS moderno definisce unità aggiuntive per un controllo più fine:

\begin{itemize}
  \item \code{svh} (\textit{small viewport height}): l'altezza del viewport
        calcolata quando i controlli sono \textbf{visibili};
  \item \code{lvh} (\textit{large viewport height}): l'altezza massima
        possibile, calcolata con i controlli \textbf{nascosti}; è il
        comportamento di \code{vh};
  \item \code{dvh} (\textit{dynamic viewport height}): l'altezza
        \textbf{corrente}, che cambia dinamicamente. Se i controlli sono
        visibili coincide con \code{svh}, cioè con \code{lvh} meno l'altezza dei
        controlli; se sono nascosti coincide con \code{lvh}.
\end{itemize}

\begin{figure}[H]
  \centering
  \begin{tikzpicture}[font=\Lato\scriptsize,
      telefono/.style={rounded corners=7pt, draw=neutral!60!black,
        line width=1.2pt, fill=white},
      barra/.style={fill=neutral!18!white, draw=neutral!45!white, line width=0.5pt},
      svh/.style={fill=purple-gradient-6!85!white, draw=purple-gradient-6!60!black},
      lvh/.style={fill=primary!60!white, draw=primary!60!black},
      dvh/.style={fill=css-deeppink!70!white, draw=css-deeppink!70!black},
      dietro/.style={fill=primary!22!white, draw=primary!60!black, dashed}]
    \foreach \x/\titolo/\visibili in {0/{controlli visibili}/1, 5.6/{controlli nascosti}/0} {
      \begin{scope}[xshift=\x cm]
        \draw[telefono] (-0.15,-0.25) rectangle (3.75,5.05);
        \node[anchor=south, font=\Lato\small\bfseries] at (1.8,5.15) {\titolo};
        % area dello schermo: da y=0 a y=4.6
        \ifnum\visibili=1
          \fill[barra] (0,0) rectangle (3.6,0.55);
          % svh e dvh: lo spazio visibile; lvh prosegue sotto la barra
          \draw[svh] (0.25,0.55) rectangle (1.05,4.6);
          \draw[dietro] (1.4,0) rectangle (2.2,0.55);
          \draw[lvh] (1.4,0.55) rectangle (2.2,4.6);
          \draw[dvh] (2.55,0.55) rectangle (3.35,4.6);
          \node[anchor=north, text=neutral!35!black] at (1.8,-0.32)
            {in basso, la barra degli indirizzi};
        \else
          \draw[svh] (0.25,0.55) rectangle (1.05,4.6);
          \draw[lvh] (1.4,0) rectangle (2.2,4.6);
          \draw[dvh] (2.55,0) rectangle (3.35,4.6);
          \node[anchor=north, text=neutral!35!black] at (1.8,-0.32)
            {la barra è nascosta};
        \fi
        \node[rotate=90, text=white, font=\Lato\scriptsize\bfseries] at (0.65,2.6) {100svh};
        \node[rotate=90, text=white, font=\Lato\scriptsize\bfseries] at (1.8,2.6) {100lvh};
        \node[rotate=90, text=white, font=\Lato\scriptsize\bfseries] at (2.95,2.6) {100dvh};
      \end{scope}
    }
    \node[anchor=west, align=left, text=neutral!25!black] at (10.0,2.4)
      {\code{svh}: sempre lo spazio\\ \phantom{\code{svh}: }con i controlli\\[4pt]
       \code{lvh}: sempre lo spazio\\ \phantom{\code{lvh}: }senza i controlli\\[4pt]
       \code{dvh}: lo spazio\\ \phantom{\code{dvh}: }di questo momento};
  \end{tikzpicture}
  \caption{Tre riquadri alti \code{100svh}, \code{100lvh} e \code{100dvh} sullo
           stesso telefono. Con la barra degli indirizzi visibile, la parte
           tratteggiata di \code{100lvh} finisce sotto la barra.}
\end{figure}

In pratica \code{100svh} garantisce che un elemento sia visibile per intero
anche quando i controlli occupano parte dello schermo, mentre \code{100dvh}
segue lo spazio effettivamente disponibile; quest'ultimo, cambiando mentre si
scorre la pagina, può però costringere il browser a ricalcolare il layout.

\section{Pattern di layout moderni}

\subsection{gap con Grid e Flexbox}

La proprietà \code{gap}, già incontrata con Grid nel capitolo 4, si usa sia sui
contenitori \textbf{grid} sia sui contenitori \textbf{flex} per aggiungere spazio
fra i figli. A differenza dei margini sui singoli figli, \code{gap} aggiunge
spazio \textbf{solo fra un elemento e l'altro}, mai fra gli elementi e il bordo
del contenitore: lo spazio sul bordo resta compito del \code{padding}.

\begin{affianco}[0.42][c]
\begin{lstlisting}[style=css, name=gap.css]
.container {
  background-color: purple;
  display: flex;
  gap: 30px;
  padding: 10px;
  flex-wrap: wrap;
}

.item {
  flex: 1 1 100px;
  height: 100px;
  background-color: orangered;
}
\end{lstlisting}
\begin{lstlisting}[style=html, name=gap.html]
<div class="container">
  <div class="item"></div>
  <div class="item"></div>
  <div class="item"></div>
  <div class="item"></div>
  <div class="item"></div>
</div>
\end{lstlisting}
\risultato
\begin{immagine}[Verso il bordo 10px di \code{padding}; fra gli elementi 30px
                 di \code{gap},\\ anche fra le righe quando gli elementi vanno
                 a capo.]
\begin{tikzpicture}[x=0.01cm, y=0.01cm, font=\Lato\scriptsize,
    it/.style={fill=css-orangered},
    misura/.style={{Stealth[length=3pt]}-{Stealth[length=3pt]}, line width=0.5pt,
      draw=neutral!25!black}]
  % contenitore largo 680px: 5 elementi da 108px
  \fill[css-purple] (0,0) rectangle (680,120);
  \foreach \i in {0,...,4} {
    \fill[it] ({10+\i*138},10) rectangle ({118+\i*138},110);
  }
  \draw[misura] (118,130) -- (148,130);
  \node[anchor=south] at (133,133) {gap: 30px};
  \draw[misura] (670,60) -- (680,60);
  \node[anchor=west] at (688,60) {padding: 10px};
  \node[anchor=north] at (340,-5) {contenitore largo 680px: una riga};
  % contenitore largo 400px: 3 + 2 elementi
  \begin{scope}[yshift=-3.25cm]
    \fill[css-purple] (0,0) rectangle (400,250);
    \foreach \i in {0,1,2} {
      \fill[it] ({10+\i*136.67},140) rectangle ({116.67+\i*136.67},240);
    }
    \fill[it] (10,10) rectangle (185,110);
    \fill[it] (215,10) rectangle (390,110);
    \draw[misura] (420,110) -- (420,140);
    \node[anchor=west] at (430,125) {gap: 30px};
    \node[anchor=north] at (200,-5) {largo 400px: due righe};
  \end{scope}
\end{tikzpicture}
\end{immagine}
\end{affianco}

Con \code{flex: 1 1 100px} ogni elemento parte da 100 pixel e si allarga per
occupare lo spazio rimasto. Nel contenitore stretto ne stanno solo tre per
riga ($3 \cdot 100 + 2 \cdot 30 = 360$ pixel sui 380 disponibili): con
\code{flex-wrap: wrap} gli altri due vanno a capo e si spartiscono la seconda
riga.

\subsection{La scorciatoia flex}

La proprietà \code{flex} è una scorciatoia (\textit{shorthand}) per tre
proprietà dei flex item, già incontrate in parte nel capitolo 4:

\begin{affianco}[0.42][c]
\begin{lstlisting}[style=css, name=flex.css]
.item {
  flex: 1 1 100px;
}

/* equivalente alla regola sopra */
.item {
  flex-grow: 1;
  flex-shrink: 1;
  flex-basis: 100px;
}
\end{lstlisting}
\risultato
\begin{itemize}
  \item \code{flex-grow}: come l'elemento \textbf{cresce}, cioè quale quota
        dello spazio libero riceve;
  \item \code{flex-shrink}: come l'elemento \textbf{si restringe} quando lo
        spazio non basta;
  \item \code{flex-basis}: la sua \textbf{dimensione di base}, lungo l'asse
        principale, prima che lo spazio venga distribuito.
\end{itemize}
\end{affianco}

Il browser parte dalle dimensioni di base di tutti gli elementi e le confronta
con lo spazio del contenitore. Se avanza spazio, lo distribuisce in proporzione
a \code{flex-grow}; se manca, toglie la differenza in proporzione a
\code{flex-shrink}. Le figure che seguono lo mostrano con tre elementi in fila;
le demo interattive di MDN, mostrate a lezione, permettono di provare dal vivo
i tre valori.

\begin{figure}[H]
  \centering
  \begin{tikzpicture}[x=0.012cm, y=0.012cm, font=\Lato\scriptsize,
      base/.style={fill=primary!35!white, draw=primary!65!black, line width=0.5pt},
      extra/.style={fill=primary!10!white, draw=primary!65!black, line width=0.5pt},
      cont/.style={draw=neutral!50!black, line width=0.6pt, dashed},
      lab/.style={anchor=east, font=\FiraCodeNL\scriptsize}]
    % contenitore di 500px, basi di 100px: avanzano 200px
    \foreach \nr/\riga/\pa/\pb/\pc in {0/{1 1 1}/66.67/66.67/66.67,
                                       1/{2 1 1}/100/50/50,
                                       2/{0 1 1}/0/100/100} {
      \begin{scope}[yshift=-\nr*0.75cm]
        \node[lab] at (-12,20) {flex-grow: \riga};
        \draw[cont] (0,0) rectangle (500,40);
        \pgfmathsetmacro{\xa}{100+\pa}
        \pgfmathsetmacro{\xb}{\xa+100+\pb}
        \draw[base] (0,0) rectangle (100,40);
        \draw[extra] (100,0) rectangle (\xa,40);
        \draw[base] (\xa,0) rectangle ({\xa+100},40);
        \draw[extra] ({\xa+100},0) rectangle (\xb,40);
        \draw[base] (\xb,0) rectangle ({\xb+100},40);
        \draw[extra] ({\xb+100},0) rectangle (500,40);
        \pgfmathtruncatemacro{\la}{round(100+\pa)}
        \pgfmathtruncatemacro{\lb}{round(100+\pb)}
        \pgfmathtruncatemacro{\lc}{round(100+\pc)}
        \node[anchor=west] at (510,20) {\la\ + \lb\ + \lc\ px};
      \end{scope}
    }
    \node[anchor=south west, font=\Lato\scriptsize] at (0,48)
      {\tikz\draw[base] (0,0) rectangle (8pt,6pt); base di 100px\quad
       \tikz\draw[extra] (0,0) rectangle (8pt,6pt); quota dei 200px liberi};
  \end{tikzpicture}
  \caption{\code{flex-grow}: in un contenitore di 500 pixel tre elementi con
           \code{flex-basis: 100px} lasciano 200 pixel liberi, divisi in
           proporzione a \code{flex-grow}. Con 2, 1, 1 il primo riceve metà dello
           spazio libero; con 0 non cresce affatto.}
\end{figure}

\begin{figure}[H]
  \centering
  \begin{tikzpicture}[x=0.012cm, y=0.012cm, font=\Lato\scriptsize,
      base/.style={fill=primary!35!white, draw=primary!65!black, line width=0.5pt},
      fuori/.style={draw=css-orangered, line width=0.6pt, dashed},
      cont/.style={draw=neutral!50!black, line width=0.6pt, dashed},
      lab/.style={anchor=east, font=\FiraCodeNL\scriptsize}]
    % basi di 200px in un contenitore di 300px: mancano 300px
    \node[lab] at (-12,20) {basi};
    \draw[cont] (0,0) rectangle (300,40);
    \foreach \i in {0,1,2} \draw[fuori] ({\i*200},4) rectangle ({\i*200+200},36);
    \node[anchor=west, text=css-orangered!80!black] at (610,20)
      {600px in 300px: ne mancano 300};
    \foreach \nr/\riga/\pa/\pb/\pc in {1/{1 1 1}/100/100/100,
                                       2/{0 1 1}/200/50/50,
                                       3/{2 1 1}/50/125/125} {
      \begin{scope}[yshift=-\nr*0.75cm]
        \node[lab] at (-12,20) {flex-shrink: \riga};
        \draw[cont] (0,0) rectangle (300,40);
        \draw[base] (0,0) rectangle (\pa,40);
        \draw[base] (\pa,0) rectangle ({\pa+\pb},40);
        \draw[base] ({\pa+\pb},0) rectangle (300,40);
        \node[anchor=west] at (310,20) {\pa\ + \pb\ + \pc\ px};
      \end{scope}
    }
  \end{tikzpicture}
  \caption{\code{flex-shrink}: tre basi di 200 pixel non entrano in un
           contenitore di 300. I 300 pixel mancanti si tolgono in proporzione a
           \code{flex-shrink} \textbf{moltiplicato per la base}: con 0 il primo
           elemento non si restringe; con 2, 1, 1 perde il doppio degli altri.}
\end{figure}

Nella restrizione il fattore viene pesato con la dimensione di base: a parità
di \code{flex-shrink}, un elemento più grande cede più spazio. Un elemento,
inoltre, non scende di norma sotto la larghezza minima del proprio contenuto
(la parola più lunga, per un testo).

\begin{figure}[H]
  \centering
  \begin{tikzpicture}[x=0.0125cm, y=0.0125cm, font=\Lato\scriptsize,
      el/.style={fill=primary!12!white, draw=primary!65!black, line width=0.5pt},
      cont/.style={draw=neutral!50!black, line width=0.6pt, dashed},
      lab/.style={anchor=east, font=\FiraCodeNL\scriptsize, align=right},
      txt/.style={font=\Lora\scriptsize, align=center}]
    \foreach \nr/\riga/\pa/\pb/\pc in {0/{flex: 1 1 auto}/127/227/127,
                                       1/{flex: 1 1 0}/160/160/160,
                                       2/{il primo: flex: 1 1 200px\\ gli altri: flex: 1 1 0}/293/93/93} {
      \begin{scope}[yshift=-\nr*1.25cm]
        \node[lab] at (-12,38) {\riga};
        \draw[cont] (0,0) rectangle (480,76);
        \draw[el] (0,0) rectangle (\pa,76);
        \draw[el] (\pa,0) rectangle ({\pa+\pb},76);
        \draw[el] ({\pa+\pb},0) rectangle (480,76);
        \pgfmathsetlengthmacro{\wa}{(\pa-8)*0.0125cm}
        \pgfmathsetlengthmacro{\wb}{(\pb-8)*0.0125cm}
        \pgfmathsetlengthmacro{\wc}{(\pc-8)*0.0125cm}
        \node[txt, text width=\wa] at ({\pa/2},38) {Uno};
        \node[txt, text width=\wb] at ({\pa+\pb/2},38) {Un testo più lungo};
        \node[txt, text width=\wc] at ({\pa+\pb+\pc/2},38) {Tre};
        \node[anchor=west] at (490,38) {\pa\ + \pb\ + \pc\ px};
      \end{scope}
    }
  \end{tikzpicture}
  \caption{\code{flex-basis} in un contenitore di 480 pixel, supponendo che i
           tre testi occupino 30, 130 e 30 pixel. Con \code{auto} la base è il
           contenuto, e lo spazio libero si aggiunge sopra; con \code{0} tutto
           lo spazio viene diviso in parti uguali; una base di 200px parte già
           in vantaggio.}
\end{figure}

\info{I valori più comuni di flex}{
  \begin{tabular}{@{}l@{\hspace{0.8cm}}l@{\hspace{0.8cm}}l@{}}
    \code{flex: 1}    & equivale a \code{1 1 0\%}  & parti uguali, indipendenti dal contenuto \\
    \code{flex: auto} & equivale a \code{1 1 auto} & parte dal contenuto, poi cresce \\
    \code{flex: none} & equivale a \code{0 0 auto} & dimensione fissa, data dal contenuto \\
    valore iniziale   & \code{0 1 auto}            & non cresce, ma può restringersi \\
  \end{tabular}}

\section{Container query}

Le media query osservano il \textbf{viewport}. Consideriamo una card che, sotto
i 300 pixel di larghezza, diventa verde e riduce il testo.

\begin{affianco}[0.5][c]
\begin{lstlisting}[style=css, name=card.css]
.card {
  background-color: lightblue;
  padding: 10px;
  border-radius: 5px;
  border: 1px solid black;
}

.card :is(h2, p) {
  font-size: 1.25em;
}

@media (max-width: 300px) {
  .card {
    background-color: lightgreen;
  }
  .card :is(h2, p) {
    font-size: 1em;
  }
}
\end{lstlisting}
\begin{lstlisting}[style=html, name=card.html]
<div class="card">
  <h2>Card title</h2>
  <p>Card text</p>
</div>
\end{lstlisting}
\risultato
\begin{browser}[finestra larga]
  \noindent\colorbox{css-lightblue}{\parbox{\dimexpr\linewidth-2\fboxsep}{%
    \vspace{2pt}{\bfseries\Large Card title}\par\vspace{6pt}
    {\large Card text}\vspace{2pt}}}
\end{browser}
\begin{minipage}{4.2cm}
\begin{browser}[sotto i 300px]
  \noindent\colorbox{css-lightgreen}{\parbox{\dimexpr\linewidth-2\fboxsep}{%
    \vspace{2pt}{\bfseries\normalsize Card title}\par\vspace{4pt}
    {\small Card text}\vspace{2pt}}}
\end{browser}
\end{minipage}
\end{affianco}

Il problema è che la stessa card può comparire in contenitori diversi: in uno
\textbf{grande}, come l'area principale del contenuto, oppure in uno
\textbf{piccolo}, come una barra laterale o una colonna stretta di una griglia.
Quando sta in un contenitore piccolo vorremmo che si comportasse come in un
viewport piccolo, \textbf{anche se il viewport è grande}. La media query non
può aiutarci: in una finestra larga la sua condizione è falsa per tutte e due
le card.

\begin{affianco}[0.44][c]
\begin{lstlisting}[style=html, name=pagina.html]
<main style="width: 350px;">
  <div class="card">
    <h2>Main card</h2>
    <p>Main text</p>
  </div>
</main>
<aside style="width: 200px;">
  <div class="card">
    <h2>Aside title</h2>
    <p>aside text</p>
  </div>
</aside>
\end{lstlisting}
\risultato
\begin{browser}[pagina.html, finestra larga]
\begin{tikzpicture}[font=\Lora,
    card/.style={draw=black, line width=0.4pt, rounded corners=2pt,
      fill=css-lightblue, anchor=north west, align=left, inner sep=5pt}]
  \node[card, text width=3.85cm] at (0,0)
    {{\bfseries\large Main card}\\[5pt] {\normalsize Main text}};
  \node[card, text width=2.05cm] at (4.35,0)
    {{\bfseries\large Aside title}\\[5pt] {\normalsize aside text}};
\end{tikzpicture}
\end{browser}
\wtnota{Il viewport è largo: entrambe le card restano azzurre.}
\end{affianco}

Le \term{container query} sono state introdotte per superare questo limite:
permettono di dare stile a un componente in base alle dimensioni di un suo
\textbf{antenato}, il \term{contenitore}.

\begin{itemize}
  \item Sull'antenato va dichiarata la proprietà \code{container-type}. Con il
        valore \code{inline-size} il contenitore si può interrogare sulla
        dimensione \textit{inline}, che per il testo orizzontale è la
        larghezza.
  \item Poi, con la regola \code{@container}, si scrivono le query: valgono per
        i \textbf{discendenti} del contenitore più vicino, non per il
        contenitore stesso.
\end{itemize}

\begin{affianco}[0.44][c]
\begin{lstlisting}[style=css, name=container.css]
aside {
  container-type: inline-size;
}

@container (max-width: 300px) {
  .card {
    background-color: lightgreen;
  }
  .card :is(h2, p) {
    font-size: 1em;
  }
}
\end{lstlisting}
\risultato
\begin{browser}[pagina.html, finestra larga]
\begin{tikzpicture}[font=\Lora,
    card/.style={draw=black, line width=0.4pt, rounded corners=2pt,
      anchor=north west, align=left, inner sep=5pt}]
  \node[card, fill=css-lightblue, text width=3.85cm] at (0,0)
    {{\bfseries\large Main card}\\[5pt] {\normalsize Main text}};
  \node[card, fill=css-lightgreen, text width=2.05cm] at (4.35,0)
    {{\bfseries\footnotesize Aside title}\\[3pt] {\footnotesize aside text}};
\end{tikzpicture}
\end{browser}
\wtnota{L'\tg{aside} è largo 200px, quindi meno di 300:\\ la sua card diventa
        verde. Il \tg{main} non è un contenitore.}
\end{affianco}

\section{Immagini responsive}

Il CSS stabilisce \textbf{quanto grande} viene mostrata un'immagine. Scaricare
un'immagine ad alta risoluzione per poi rimpicciolirla su uno schermo piccolo,
però, non è efficiente: si trasferiscono byte che nessuno vedrà. L'HTML permette
di fornire \textbf{più file} per la stessa immagine, i \term{candidati}, e
lascia al browser la scelta di quello più adatto al dispositivo. Il vantaggio è
che il browser scarica \textbf{solo} la versione che gli serve. Si usano
l'attributo \attr{srcset} oppure l'elemento \tg{picture}.

\subsection{srcset}

\begin{lstlisting}[style=html, name=srcset.html]
<img src="fu-tzu-400.jpg"
     srcset="fu-tzu-400.jpg 400w,
             fu-tzu-800.jpg 800w,
             fu-tzu-1200.jpg 1200w"
     alt="Portrait of the great master Fu-Tzu">
\end{lstlisting}

All'immagine sono associati tre file diversi.

\begin{itemize}
  \item L'attributo \attr{src} è il valore di riserva, per i browser che non
        supportano \attr{srcset}.
  \item \attr{srcset} è un elenco di nomi di file, ciascuno seguito da un
        \textbf{descrittore}:
        \begin{itemize}
          \item i \textbf{descrittori di larghezza}, un numero seguito da
                \code{w}, indicano la larghezza intrinseca del file in pixel;
          \item i \textbf{descrittori di densità}, un numero seguito da
                \code{x}, indicano per quale densità di pixel del dispositivo è
                pensato il file, rispetto alla densità standard. Per esempio
                \code{srcset="logo.png 1x, logo-2x.png 2x"} serve il secondo
                file agli schermi ad alta densità.
        \end{itemize}
\end{itemize}

Con i descrittori di larghezza il browser deve sapere quanto sarà larga
l'immagine nella pagina: lo indica l'attributo \attr{sizes}, che se manca vale
\code{100vw}, cioè tutta la larghezza del viewport. Moltiplicando quella
larghezza per la densità dello schermo ottiene i pixel necessari, e sceglie il
candidato più adatto. Per esempio, con l'immagine larga quanto il viewport:

\begin{center}
\renewcommand{\arraystretch}{1.35}
\begin{tabular}{@{}l@{\hspace{0.9cm}}c@{\hspace{0.9cm}}c@{\hspace{0.9cm}}c@{\hspace{0.9cm}}l@{}}
  \textbf{Dispositivo} & \textbf{Viewport} & \textbf{Densità} &
    \textbf{Pixel necessari} & \textbf{File scaricato} \\
  smartphone     & 375px  & 2 & 750  & \code{fu-tzu-800.jpg} \\
  tablet         & 768px  & 1 & 768  & \code{fu-tzu-800.jpg} \\
  portatile      & 1440px & 1 & 1440 & \code{fu-tzu-1200.jpg} (il più grande) \\
\end{tabular}
\end{center}

La scelta finale spetta comunque al browser, che può tenere conto anche della
cache o della qualità della connessione.

\subsection{picture e art direction}

\begin{affianco}[0.6][c]
\codicepiccolo
\begin{lstlisting}[style=html, name=picture.html]
<picture>
  <source media="(max-width: 200px)" srcset="fu-tzu-400.png">
  <source media="(max-width: 400px)" srcset="fu-tzu-800.png">
  <img src="fu-tzu-1200.png" alt="Portrait of the great master Fu-Tzu">
</picture>
\end{lstlisting}
\risultato
\begin{itemize}
  \item il browser usa il primo \tg{source} la cui media query è vera;
  \item se nessuna lo è, usa l'\tg{img};
  \item l'\tg{img} è obbligatoria: è l'elemento che viene davvero
        mostrato, i \tg{source} ne cambiano solo il file.
\end{itemize}
\end{affianco}

L'\term{art direction} è il processo di scelta del contenuto più adatto a
specifiche condizioni. Non si tratta solo di versioni a risoluzioni diverse
della stessa immagine: su uno smartphone si può mostrare, per esempio, un
grafico con meno dettagli, o una foto ritagliata sul soggetto, e su un desktop
la versione completa.

\begin{figure}[H]
\centering
\begin{minipage}[t]{4.6cm}\vspace{0pt}
\begin{browser}[picture.html]
  \centering
  \begin{tikzpicture}
    \node[inner sep=0pt] (img) {\includegraphics[width=3.7cm]{sources/fu-tzu.jpg}};
    \node[anchor=north, text=white, font=\Lato\bfseries\Huge] at ([yshift=-2pt]img.north) {1200};
  \end{tikzpicture}
\end{browser}
\end{minipage}\hfill
\begin{minipage}[t]{14.2cm}\vspace{0pt}
\begin{devtools}{Rete}
  {\Lato\scriptsize\setlength{\tabcolsep}{5pt}%
  \colorbox{dtbar}{\makebox[\dimexpr\linewidth-2\fboxsep][l]{%
    \color{neutral!45!black}\quad fu-tzu\hfill}}\par
  \noindent\hspace*{6pt}{\color{neutral!30!black}Tutti\quad HTML\quad CSS\quad
    JS\quad XHR\quad Caratteri\quad \colorbox{dtline}{Immagini}\quad Media\quad
    WS\quad Altro}\par\vspace{2pt}
  {\color{dtline}\hrule}
  \noindent\begin{tabular*}{\linewidth}{@{\hspace{8pt}}l@{\extracolsep{\fill}}llllll@{\hspace{8pt}}}
    \rule{0pt}{10pt}\textbf{Stato} & \textbf{Metodo} & \textbf{Dominio} &
    \textbf{File} & \textbf{Iniziatore} & \textbf{Tipo} & \textbf{Dimensione} \\ \hline
    \rule{0pt}{10pt}\dtstato{200} & GET & 127.0.0.1:3004 & fu-tzu-1200.png &
      imageset & png & 1,65 MB \\
  \end{tabular*}\par\vspace{8pt}
  {\color{dtline}\hrule}\vspace{2pt}
  \noindent\hspace*{6pt}{\color{neutral!25!black}1 richiesta (filtrata)}\par\vspace{3pt}}
\end{devtools}
\end{minipage}
\caption{In una finestra larga più di 400 pixel nessun \tg{source} è adatto:
         la scheda \textit{Rete}, filtrata sulle richieste di immagini
         \code{fu-tzu}, mostra che il browser ha scaricato solo
         \code{fu-tzu-1200.png}, non le altre due versioni.}
\end{figure}

\section{Media query sulle preferenze dell'utente}

Le media query possono verificare molto più delle dimensioni del viewport.
Possono riflettere anche le \textbf{preferenze dell'utente}, che di solito
vengono dalle impostazioni del browser o del sistema operativo.

\begin{center}
\renewcommand{\arraystretch}{1.4}
\begin{tabular}{@{}l@{\hspace{1.0cm}}l@{}}
  \textbf{Preferenza} & \textbf{Media query} \\
  schema di colori scuro          & \code{@media (prefers-color-scheme: dark)} \\
  meno animazioni                 & \code{@media (prefers-reduced-motion: reduce)} \\
  più (o meno) contrasto          & \code{@media (prefers-contrast: more)} \\
  meno trasparenze                & \code{@media (prefers-reduced-transparency: reduce)} \\
\end{tabular}
\end{center}

Nelle slide il valore delle due preferenze \textit{reduced} è scritto
\code{reduced}: quello corretto è \code{reduce}, e con un valore sbagliato la
media query viene ignorata.

\begin{affianco}[0.46][c]
\begin{lstlisting}[style=css, name=colori.css]
:root {
  --bg-color: white;
  --text-color: indigo;
}

body {
  background-color: var(--bg-color);
  color: var(--text-color);
}

@media (prefers-color-scheme: dark) {
  :root {
    --bg-color: indigo;
    --text-color: white;
  }
}
\end{lstlisting}
\risultato
\begin{center}
\begin{tikzpicture}[font=\Lato\scriptsize,
    opz/.style={draw=neutral!45!white, line width=0.6pt, rounded corners=3pt,
      minimum width=1.6cm, minimum height=0.7cm, fill=white}]
  \node[anchor=south, font=\Lato\scriptsize\bfseries] at (2.0,0.55)
    {impostazioni del browser: aspetto dei siti web};
  \node[opz] (s) at (0,0) {sistema};
  \node[opz, draw=dtblue, line width=1pt] (c) at (2.0,0) {chiaro};
  \node[opz, fill=neutral!75!black, text=white] (d) at (4.0,0) {scuro};
\end{tikzpicture}
\end{center}
\begin{minipage}[t]{0.49\linewidth}
\begin{browser}[chiaro]
  {\bfseries\large\textcolor{css-indigo}{Color Scheme}}\par\vspace{3pt}
  {\scriptsize\textcolor{css-indigo}{Color scheme preference.}}
\end{browser}
\end{minipage}\hfill
\begin{minipage}[t]{0.49\linewidth}
\begin{browser}[scuro]
  % lo sfondo indaco copre anche i margini interni della finestra
  \vspace*{-6pt}\noindent\makebox[\linewidth][l]{\hspace*{-8pt}%
  \colorbox{css-indigo}{\parbox{\dimexpr\linewidth+16pt-2\fboxsep}{%
    \vspace{6pt}\hspace*{4pt}{\bfseries\large\textcolor{white}{Color Scheme}}\par\vspace{3pt}
    \hspace*{4pt}{\scriptsize\textcolor{white}{Color scheme preference.}}\vspace{6pt}}}}%
  \vspace*{-6pt}
\end{browser}
\end{minipage}
\wtnota{Con \textit{sistema} il browser segue l'impostazione\\ del sistema
        operativo.}
\end{affianco}

Le custom property si prestano bene a questo uso: la media query ridefinisce
solo le variabili, e tutte le regole che le usano cambiano di conseguenza.

\section{@supports}

La regola \code{@supports} applica degli stili solo quando il browser supporta
una certa proprietà con un certo valore.

\begin{affianco}[0.46][c]
\begin{lstlisting}[style=css, name=supports.css]
.cards {
  display: block;
}

@supports (display: grid) {
  .cards {
    display: grid;
    grid-template-columns: 1fr 1fr;
  }
}
\end{lstlisting}
\risultato
\begin{immagine}[Lo stesso contenuto, senza e con il supporto per Grid.]
\begin{tikzpicture}[font=\Lato\scriptsize,
    c/.style={draw=primary!60!black, fill=primary!12!white, line width=0.6pt,
      minimum height=0.5cm, inner sep=2pt}]
  \node[anchor=south] at (1.4,1.85) {senza Grid};
  \foreach \i in {0,...,3} \node[c, minimum width=2.8cm] at (1.4,{1.5-\i*0.6}) {card \the\numexpr\i+1\relax};
  \node[anchor=south] at (5.2,1.85) {con Grid};
  \foreach \i in {0,1} {
    \node[c, minimum width=1.6cm] at (4.55,{1.5-\i*0.6}) {card \the\numexpr2*\i+1\relax};
    \node[c, minimum width=1.6cm] at (5.85,{1.5-\i*0.6}) {card \the\numexpr2*\i+2\relax};
  }
\end{tikzpicture}
\end{immagine}
\end{affianco}

La prima regola è il valore di riserva (\textit{fallback}): un browser che non
conosce Grid mostra le card una sotto l'altra, che è comunque leggibile. Il
layout avanzato su due colonne viene applicato solo dai browser che lo
supportano. È un esempio di \term{progressive enhancement}: si parte da una
base che funziona ovunque e la si arricchisce dove possibile.

\section{Annidamento nativo}

Nel CSS moderno le regole dei selettori correlati si possono
\textbf{annidare}, come Sass permetteva di fare da anni (capitolo 19): una
regola scritta dentro un'altra si applica ai \textbf{discendenti} degli elementi
selezionati da quella esterna.

\begin{affianco}[0.5][c]
\begin{lstlisting}[style=css, name=annidato.css]
.card {
  background-color: lightblue;
  padding: 10px;
  border-radius: 5px;
  border: 1px solid black;

  :is(h2, p) {
    font-size: 1.25em;
  }
}
\end{lstlisting}
\risultato
\begin{lstlisting}[style=css, name=equivalente.css]
.card {
  background-color: lightblue;
  padding: 10px;
  border-radius: 5px;
  border: 1px solid black;
}

.card :is(h2, p) {
  font-size: 1.25em;
}
\end{lstlisting}
\wtnota{A sinistra la regola annidata, a destra la sua forma equivalente.}
\end{affianco}

\subsection{Il selettore \&}

Nei selettori annidati il simbolo \code{\&} si riferisce esplicitamente al
selettore della regola genitore. Serve quando la regola annidata non riguarda
un discendente, ma l'elemento \textbf{stesso}, per esempio con una
pseudo-classe, oppure quando si vuole un combinatore diverso dallo spazio.

\begin{affianco}[0.5][c]
\begin{lstlisting}[style=css, name=nesting.css]
.card {
  background: white;

  &:hover {
    background: lavender;
  }

  & > h2 {
    color: rebeccapurple;
  }
}
\end{lstlisting}
\risultato
\begin{lstlisting}[style=css, name=equivalente.css]
.card {
  background: white;
}

.card:hover {
  background: lavender;
}

.card > h2 {
  color: rebeccapurple;
}
\end{lstlisting}
\end{affianco}

Senza \code{\&}, un selettore annidato viene combinato con quello esterno come
\textbf{discendente}: \code{:hover} annidato dentro \code{.card} equivarrebbe a
\code{.card :hover}, cioè a qualunque discendente della card sotto il
puntatore, e non alla card stessa. È la differenza fra \code{.card:hover} e
\code{.card :hover}, che sta tutta in uno spazio.

\warning{Annidare con moderazione}{
  Conviene mantenere l'annidamento \textbf{poco profondo}. Un annidamento
  eccessivo rende più difficili da capire sia i selettori che ne risultano sia
  il comportamento della cascata: con quattro o cinque livelli diventa difficile
  ricostruire a mente quale selettore completo si sta scrivendo, e quale sia la
  sua specificità.}

\section{Riferimenti}

\begin{itemize}
  \item \textbf{Learn CSS} --- web.dev. \\
        \urlx{https://web.dev/learn/css/} \\
        Sezioni 14 e 17.
  \item \textbf{Using CSS custom properties (variables)} --- MDN web docs. \\
        \urlx{https://developer.mozilla.org/en-US/docs/Web/CSS/Guides/Cascading_variables/Using_custom_properties}
  \item \textbf{calc() CSS function} --- MDN web docs. \\
        \urlx{https://developer.mozilla.org/en-US/docs/Web/CSS/Reference/Values/calc}
\end{itemize}
